\section{Tele-Odontologia e as Redes Descentralizadas de Cuidado}
\label{sec:teleodontologia}

A tele-odontologia, rigorosamente demarcada como a sinergia entre Tecnologias da Informação e Comunicação (TIC) e a práxis estomatológica remota, alçou um status de proeminência incontornável após a reestruturação dos sistemas de saúde globais iniciada em 2020 \cite{ijmi2026disability}. Distanciando-se de um mero artifício de videoconferência, ela encapsula protocolos de teletriagem assíncrona, telemonitoramento cinemático, teleinterconsulta especializada e teleorientação, estabelecendo um \textit{continuum} de cuidados preventivos e curativos.

No arcabouço jurisprudencial brasileiro, a Resolução CFO 226/2020 forneceu guarida regulatória a essa prática, consolidando modalidades não-invasivas à distância. Embora a legislação sabiamente vede o estabelecimento de prognósticos invasivos e intervenções reabilitadoras extramuros sem um exame clínico físico validado primariamente, as possibilidades abertas para a odontologia em saúde pública são substanciais. A tele-odontologia está profundamente alicerçada às prerrogativas normativas e programáticas do Ministério da Saúde e da Organização Mundial da Saúde (OMS), servindo de baluarte para a democratização da promoção de saúde.

\destaque{A virtualização dos canais de cuidado revela sua máxima potência mitigadora nas iniquidades sociogeográficas históricas que afligem os ecossistemas de saúde globais \cite{du2020marginalized}.}

A literatura sublinha que a ausência ou má-distribuição de mão de obra odontológica hiperespecializada induz a uma morbidade desproporcional. Estudo empírico proeminente de Frota e colaboradores demonstra inequivocamente que a agregação ostensiva da tele-odontologia à capilaridade do Sistema Único de Saúde (SUS) logrou otimizar o percentual de rastreio precoce e teleassistência em populações de alta vulnerabilidade, reduzindo a supressão de demanda reprimida nas extremidades territoriais norte e nordeste do Brasil \cite{frota2025sus, cuadros2023access}.

Sobretudo no panorama discutido nesta obra, a simbiose epistemológica entre a tele-odontologia e a arquitetura dos \textit{Digital Twins} apresenta-se como a evolução natural desse paradigma \cite{olawade2026advancing}. A materialização de um avatar digital hospedado na nuvem permite que o paciente seja avaliado por juntas odontológicas interdisciplinares em centros de referência terciários a milhares de quilômetros de distância. Estas plataformas suportam a análise cooperativa (\textit{crowdsourcing} acadêmico) de laudos tomográficos ou escaneamentos patológicos em tempo real, fornecendo os alicerces teóricos para as abordagens comunitárias e sistêmicas da saúde, tópico que reverbera com profusão no \autoref{ch:introducao} \cite{shen2024effectiveness}.
